% API仕様
\section{API仕様}

\subsection{IModule<T>テンプレートインターフェース}

\begin{apibox}{IModule<T> — モジュールテンプレートインターフェース}
\textbf{定義ファイル:} \texttt{lib/ModuleCore/IModule.h}

全ハードウェアモジュールが実装すべき抽象インターフェース。
C++テンプレートにより、プロジェクト固有の型（\texttt{SystemData}等）に一切依存しない。
\texttt{T}にはプロジェクトで定義した共有データ構造体の型を指定する。
\end{apibox}

\begin{lstlisting}[style=cpp, caption={IModule<T>テンプレートインターフェース定義}]
// IModule.h — プロジェクト非依存のコアヘッダー
#pragma once

// <typename T> と書くことで、
// 「Tという型（構造体）を後から指定する」という宣言になる
template <typename T>
class IModule {
public:
    // 仮想デストラクタ
    virtual ~IModule() {}

    // モジュールの初期化
    // 戻り値: true=成功, false=失敗
    virtual bool init() = 0;

    // モジュールの更新（毎ループ呼び出される）
    // 引数: data - プロジェクト共有データへの参照
    virtual void update(T& data) = 0;
};
\end{lstlisting}

\begin{notebox}{テンプレートを使う理由}
テンプレートを使わない場合、\texttt{IModule.h}内で\texttt{SystemData}を\texttt{\#include}する必要がある。
しかし\texttt{SystemData}はプロジェクト固有の型であるため、コアライブラリがプロジェクトに依存してしまう。\\
\\
テンプレートにより、\texttt{IModule.h}は何も\texttt{\#include}せずに済み、
どんなプロジェクトにもそのままコピーして流用できる。
\end{notebox}

\subsubsection{init()}

\begin{tabularx}{\textwidth}{l X}
\toprule
\textbf{項目} & \textbf{内容} \\
\midrule
戻り値 & \texttt{bool} — \texttt{true}で初期化成功、\texttt{false}で失敗 \\
呼び出しタイミング & \texttt{setup()}内で1回のみ \\
責務 & ハードウェアの初期化、ピン設定、ModuleTimerの\texttt{setTime()}呼び出し等 \\
エラー時 & \texttt{false}を返す。呼び出し側でモジュールの無効化やログ出力を行う \\
\bottomrule
\end{tabularx}

\subsubsection{update(T\& data)}

\begin{tabularx}{\textwidth}{l X}
\toprule
\textbf{項目} & \textbf{内容} \\
\midrule
引数 & \texttt{T\&} — プロジェクト共有データへの参照（nullにならない） \\
戻り値 & \texttt{void} \\
呼び出しタイミング & \texttt{loop()}内で毎ループ呼び出される \\
責務 & センサー読み取り、出力制御、共有データの更新 \\
エラー時 & 共有データ内のサブ構造体にエラーフラグを設定 \\
\bottomrule
\end{tabularx}

\subsubsection{使用例}

\begin{lstlisting}[style=cpp, caption={モジュール実装の例}]
// LedModule.h
#pragma once
#include "IModule.h"
#include "SystemData.h"  // プロジェクト固有

struct LedData {
    bool state;
    uint8_t brightness;
};

// テンプレート引数にSystemDataを指定して継承
class LedModule : public IModule<SystemData> {
public:
    bool init() override {
        pinMode(LED_PIN, OUTPUT);
        return true;
    }

    // data はポインタではなく参照なので、ドット演算子でアクセス
    void update(SystemData& data) override {
        digitalWrite(LED_PIN, data.led.state);
    }
};
\end{lstlisting}

% -------------------------------------------

\subsection{ModuleTimerクラス}

\begin{apibox}{ModuleTimer — ノンブロッキングタイマー}
\textbf{定義ファイル:} \texttt{lib/ModuleCore/ModuleTimer.h}

\texttt{millis()}ベースのシンプルなタイマークラス。
\texttt{delay()}を使用せずに周期実行・デバウンス・ワンショット・タイムアウト等を実現する。
インスタンスを複数作成することで、モジュール内で複数のタイマーを独立に管理できる。
IModuleと同様にプロジェクト非依存であり、そのまま流用可能。
\end{apibox}

\begin{lstlisting}[style=cpp, caption={ModuleTimerクラス定義}]
// ModuleTimer.h — プロジェクト非依存のコアヘッダー
#pragma once
#include <Arduino.h>

class ModuleTimer {
private:
    unsigned long _startTime;

public:
    // コンストラクタ
    ModuleTimer() : _startTime(0) {}

    // 基準時刻をセット
    // 引数なし: getNowTime()が0から開始
    // 引数あり: getNowTime()が指定値から開始
    void setTime(unsigned long offsetMs = 0) {
        _startTime = millis() - offsetMs;
    }

    // 基準時刻からの経過時間(ms)を取得
    unsigned long getNowTime() const {
        return millis() - _startTime;
    }
};
\end{lstlisting}

\subsubsection{setTime(unsigned long offsetMs = 0)}

\begin{tabularx}{\textwidth}{l X}
\toprule
\textbf{項目} & \textbf{内容} \\
\midrule
引数 & \texttt{offsetMs}（デフォルト: 0） — 初期オフセット値(ms) \\
戻り値 & \texttt{void} \\
動作 & 内部の基準時刻を\texttt{millis() - offsetMs}に設定する。
       \texttt{offsetMs = 0}の場合、直後の\texttt{getNowTime()}は0を返す。
       \texttt{offsetMs = 100}の場合、直後の\texttt{getNowTime()}は100を返す \\
用途 & タイマーの開始・リセット。\texttt{init()}内および周期リセット時に使用 \\
\bottomrule
\end{tabularx}

\subsubsection{getNowTime() const}

\begin{tabularx}{\textwidth}{l X}
\toprule
\textbf{項目} & \textbf{内容} \\
\midrule
引数 & なし \\
戻り値 & \texttt{unsigned long} — 最後の\texttt{setTime()}からの経過時間(ms) \\
動作 & \texttt{millis() - \_startTime}を返す \\
用途 & 経過時間の判定。条件分岐と組み合わせて周期実行・デバウンス等を実現 \\
\bottomrule
\end{tabularx}

\subsubsection{使用パターン}

\begin{lstlisting}[style=cpp, caption={周期実行パターン}]
ModuleTimer readTimer;

bool init() override {
    readTimer.setTime();
    return true;
}

void update(SystemData& data) override {
    if (readTimer.getNowTime() >= 1000) {
        readTimer.setTime();  // リセット
        data.temp.temperature = readSensor();
    }
}
\end{lstlisting}

\begin{lstlisting}[style=cpp, caption={デバウンスパターン}]
ModuleTimer debounce;

void update(SystemData& data) override {
    bool rawInput = digitalRead(BUTTON_PIN);

    if (rawInput != lastRawInput) {
        debounce.setTime();  // 入力変化でリセット
        lastRawInput = rawInput;
    }

    if (debounce.getNowTime() >= 50) {
        data.button.pressed = lastRawInput;  // 50ms安定で確定
    }
}
\end{lstlisting}

% -------------------------------------------

\subsection{SystemData構造体}

\begin{apibox}{SystemData — プロジェクト共有データ}
\textbf{定義ファイル:} \texttt{include/SystemData.h}（プロジェクト固有）

各モジュールが定義するサブ構造体を集約した、プロジェクト固有のデータ構造体。
全モジュールの\texttt{update()}メソッドに参照として渡される。
\texttt{IModule<T>}のテンプレート引数\texttt{T}として指定する型である。
\end{apibox}

\subsubsection{設計ルール}

\begin{enumerate}
    \item 各モジュールは、自身が管理するデータをサブ構造体として定義する（モジュールのヘッダーファイル内）
    \item \texttt{SystemData}構造体は、プロジェクトで使用するモジュールのサブ構造体をメンバとして持つ
    \item モジュールは\texttt{update()}内で、自身のサブ構造体に書き込み、必要に応じて他モジュールのサブ構造体を読み取る
\end{enumerate}

\begin{lstlisting}[style=cpp, caption={SystemData構成例}]
// SystemData.h
#pragma once
#include "LedModule.h"
#include "TempSensorModule.h"
#include "CameraModule.h"

struct SystemData {
    LedData led;
    TempData temp;
    CameraData camera;
};
\end{lstlisting}
